\section{RACE-008: Async Initialization Race in query()}
\label{sec:race-008}
\subsection*{Metadata}
\begin{tabular}{ll}
\textbf{Issue ID} & RACE-008 \\
\textbf{Severity} & MEDIUM \\
\textbf{Category} & Race Condition \\
\textbf{Branch} & \branchref{fix/RACE-008-async-init-session}
\textbf{Changes} & \branchdiff{fix/RACE-008-async-init-session} \\
\end{tabular}

\subsection*{Executive Summary}
In \srcfile{src/sdk.ts}, the \code{query()} function starts an async IIFE immediately, but the \code{\_sessionId} variable is only set after \code{loadAgent()} completes (about 100-500ms later). The \code{sessionId()} accessor returns \code{\_sessionId} synchronously. If a consumer calls \code{sessionId()} before the agent finishes loading — in a hook's \code{onSessionStart} handler or for telemetry tagging — the returned value is an empty string, breaking session correlation.

\subsection*{Root Cause Analysis}
The \code{\_sessionId} is set asynchronously but accessed synchronously:

\begin{lstlisting}[language=TypeScript]
// BEFORE — synchronous sessionId() can return ""
export function query(options: QueryOptions): Query {
    let _sessionId = options.sessionId ?? "";  // Initialized empty

    const runPromise = (async () => {
        const loaded = await loadAgent(dir, ...);  // ~100-500ms
        _sessionId = _sessionId || loaded.sessionId;  // SET here
    })();

    return {
        sessionId() {
            return _sessionId;  // MAY BE EMPTY if called before loadAgent resolves
        },
    };
}
\end{lstlisting}

\subsection*{Fix Implementation}
A promise-based initialization pattern ensures \code{sessionId()} waits for the agent to load:

\begin{lstlisting}[language=TypeScript]
// AFTER — sessionId() returns Promise<string> and waits for init
export function query(options: QueryOptions): Query {
    let _sessionId = options.sessionId ?? "";
    let _sessionIdResolve: (() => void) | null = null;
    const _sessionIdPromise = new Promise<void>((r) => { _sessionIdResolve = r; });
    if (_sessionId) _sessionIdResolve?.();

    const runPromise = (async () => {
        const loaded = await loadAgent(dir, ...);
        _sessionId = _sessionId || loaded.sessionId;
        _sessionIdResolve?.();  // Signal that sessionId is ready
    })();

    return {
        sessionId(): string | Promise<string> {
            return _sessionIdPromise.then(() => _sessionId);
        },
    };
}
\end{lstlisting}

The type signature is also updated to allow returning \code{Promise<string>}:

\begin{lstlisting}[language=TypeScript]
// src/sdk-types.ts
export interface Query extends AsyncGenerator<GCMessage, void, undefined> {
    sessionId(): string | Promise<string>;
}
\end{lstlisting}

\subsection*{Affected Files}
\begin{itemize}
    \item \srcfile{src/sdk.ts} — query() function
    \item \srcfile{src/sdk-types.ts} — Query interface
    \item \srcfile{src/\_\_tests\_\_/async-init-session.test.ts}
\end{itemize}

\subsection*{Verification}
Tests verify sessionId resolves to the correct value after initialization:

\begin{lstlisting}[language=TypeScript]
async function getSessionId(): Promise<string> {
    await _sessionIdPromise;
    return _sessionId;
}
setTimeout(() => { _sessionId = "session-123"; _sessionIdResolve!(); }, 10);
const id = await getSessionId();
assert.equal(id, "session-123");
\end{lstlisting}
